
\chapter*{Resumen}
\addcontentsline{toc}{chapter}{Resumen}

La capacidad de un analista de ciberseguridad para detectar, investigar y responder ante un incidente depende en gran medida de su comprensión del comportamiento del adversario. Sin embargo, los recursos formativos disponibles suelen centrarse en técnicas aisladas o en escenarios parciales que no reflejan la complejidad de una intrusión real, dificultando la adquisición de una visión integral del ciclo de ataque.

Este trabajo aborda dicha carencia mediante el diseño, desarrollo y ejecución controlada de un ciberataque realista en un entorno de laboratorio. El escenario cubre el ciclo completo de intrusión: acceso inicial mediante dispositivo BadUSB, evasión de AMSI, desarrollo de shellcode en ensamblador (reverse shell), inyección de código mediante DLL Hijacking (Side Loading, DLL Proxying y Reflective Loading), despliegue de un implante de Sliver C2 cargado íntegramente en memoria y establecimiento de persistencia mediante tareas programadas con LOLBins.

El resultado principal es un conjunto de artefactos forenses extraídos de los sistemas comprometidos, acompañado de una documentación completa que establece la trazabilidad entre cada acción ofensiva, la técnica MITRE ATT\&CK asociada y las evidencias donde puede observarse su rastro. Este recurso permite al analista reconstruir el incidente desde la perspectiva forense, pero también comprender las decisiones y motivaciones del atacante en cada fase.

La contribución fundamental del trabajo es proporcionar a la comunidad académica y profesional un entorno de entrenamiento que integra la perspectiva ofensiva y defensiva, reduciendo la brecha entre la teoría y la práctica en la formación de analistas de ciberseguridad.

\textbf{Palabras clave:} simulación de adversarios, red team, DFIR, análisis forense, MITRE ATT\&CK, desarrollo de malware, shellcode, DLL hijacking, reflective loading, Sliver C2, formación en ciberseguridad.